\begingroup
\let\clearpage\relax
\let\cleardoublepage\relax
\chapter{Introduction}
\endgroup

\section{Motivation}

Unexpected changes in ordinary traffic patterns can occur at any instant and may result in network faults or quality-of-service problems.
Telecom traffic is not a single univariate series; it can be viewed under several categorical dimensions.
For example, a call can have different types, an event type may take many values, and a subscriber may belong to a given type.
Traffic type, service type, partner-related information, and related categories create many candidate aggregations on which anomaly detection can be run.
Expert labels for real incidents are scarce, so candidate aggregations cannot be ranked with standard supervised metrics.
This dissertation focuses on injecting synthetic anomalies into each candidate aggregation and using \ac{AUC} on Matrix Profile discord scores as a label-free fitness signal.

Having a fully aggregated view can hide local anomalies, while monitoring every categorical combination separately is costly computationally and in terms of interpretability.
Moreover, some fully disaggregated series contain so few events that anomalies cannot be identified with confidence.
This creates a dilemma: which aggregation level should anomaly detection work on?
Operators may use a fixed hierarchy or manual rules; for example, aggregating by protocol or subscriber type can make some anomalies visible.
Running anomaly detection independently on every variable combination, however, remains expensive and hard to interpret.
Therefore, automatic search methods are needed to identify suitable candidate aggregations.

\section{Objectives}

This dissertation aims to develop a framework for anomaly detection in grouped telecommunications time series with hierarchical aggregation, using Grammatical Evolution to evolve programmes from a formal grammar.
From a data perspective, the aim is to work with the categorical dimensions of real mobile-network traffic records---call type, event type, subscriber type, traffic type, service type, partner-related information, and related fields---so that anomalies can be studied at global and local aggregation levels.

From a methodological point of view, aggregation choice is framed as a combinatorial search over those categorical dimensions.
Candidate hierarchies are encoded in a \ac{BNF} grammar so that only valid aggregation structures are generated.
Grammatical Evolution (PonyGE2) then evolves aggregation programmes from that grammar.
For each candidate, an aggregated series is built and scored with Matrix Profile (STUMPY) discord distances in a label-free way.
Because anomaly labels are scarce, synthetic anomalies provide the evaluation signal: \ac{AUC} on the discord scores is used as fitness, and evolved hierarchies are compared with fixed aggregation baselines.

The work is guided by three research questions, which are answered in Chapter~4:
\begin{enumerate}
    \item Can the choice of aggregation be framed as a combinatorial search over a grammar, so that every candidate is a valid hierarchy?
    \item Can a label-free Matrix Profile scorer, evaluated on synthetic anomalies, drive this search?
    \item Do the evolved hierarchies detect anomalies better than fixed aggregations, such as the fully aggregated series, single-feature splits and full disaggregation?
\end{enumerate}

\section{Contributions}

The main contributions of this dissertation are:
\begin{itemize}
    \item the formulation of aggregation choice for anomaly detection in grouped telecommunications series as a grammar-based search, with a fitness built from a label-free detector and synthetic anomalies;
    \item an implementation in Python on top of PonyGE2 and STUMPY that evaluates candidate hierarchies on real operator data and can be reproduced from the saved runs;
    \item an empirical comparison of evolved and fixed hierarchies on a mobile network switch. The evolved hierarchies reached a test \ac{AUC} of $0.678 \pm 0.066$ and were significantly better than the fully aggregated series and the Call, Event and Service Type splits, but not better than the best single-feature splits;
    \item evidence that conclusions about the best aggregation depend strongly on how false alarms are counted, since the event-weighted fitness ignores leaves without anomalies and can be exploited by the search.
\end{itemize}

\section{Document Organization}

The dissertation is organised into five chapters.
Chapter~1 presents the motivation, objectives, research questions, contributions and document organisation.
Chapter~2 is the background: the bibliographic search, the theory, and the literature review.
Chapter~3 describes the telecommunications dataset, preprocessing and aggregation, synthetic anomaly injection, hierarchy construction, and the PonyGE2 implementation.
Chapter~4 presents the experimental results, comparing aggregation strategies and examining the evolved hierarchical structures.
Chapter~5 summarises the contributions and limitations of the present work, and outlines future work, including the use of real incident labels and online detection.
